\label{sec:prediction}

The central prediction of the O-series is:

\begin{quote}
\textbf{The Compactness-Outcomes Hypothesis}: Compact algorithmic congressional redistricting
plans, as measured by district Polsby-Popper scores, predict better democratic outcomes
on all four O-series dimensions (competitiveness, turnout, polarization, constituent
distance), independent of the partisan composition of the state's congressional delegation.
\end{quote}

This prediction has three components that require separate justification.

\subsection{Independence from Partisan Composition}

The prediction is qualified by ``independent of partisan composition'' because one
might worry that compact plans predict better outcomes merely because they happen
to produce more Democratic-leaning delegations (given geographic sorting, which tends
to concentrate Democrats in compact urban cores). If compact plans were just
another way of saying ``plans that favor Democrats,'' then the outcome improvements
would reflect partisan selection, not geometric effects.

We address this in two ways. First, the O-series controls throughout for the partisan
composition of the state and the district under analysis. The Polsby-Popper effect is
estimated within-partisan-composition bins, so any residual compactness effect is
not merely proxying for partisan advantages. Second, the O.4 analysis provides a
particularly clean test: the mean constituent driving time outcome is not mechanically
related to partisan composition, so any Polsby-Popper correlation with driving time
cannot be explained by partisan selection.

\subsection{The Compactness Mechanism}

The theoretical mechanism connecting compactness to outcomes runs through district
geometry's effect on the political heterogeneity of district populations. Consider
a district drawn as a narrow corridor connecting two geographically distant communities
(a characteristic shape in documented gerrymanders). Such a district:

\begin{enumerate}
\item Has low Polsby-Popper because the corridor creates a large perimeter-to-area ratio.
\item Combines communities with different political preferences (the corridor may cross
  from a Democratic urban neighborhood to a Republican exurb), but places them in a
  single district where only one representative can serve both.
\item Creates representational distance: the representative's office is located somewhere
  along the corridor, maximizing distance to one end or the other.
\item Is typically designed to create a safe seat for one party (by controlling the ratio
  of urban to exurban population), reducing electoral competitiveness.
\end{enumerate}

Compact districts avoid this configuration. A district drawn to minimize perimeter
relative to area necessarily groups geographically proximate communities, which tend
to be more politically homogeneous. The political homogeneity effect cuts both ways:
very compact urban districts are safely Democratic; very compact rural districts are
safely Republican. But the marginal compact districts---those at the edge of geographic
sorting---are more competitive than the marginal non-compact districts precisely because
they have not been engineered to achieve a target partisan split.

\subsection{Quantitative Predictions}

Table~\ref{tab:predictions} summarizes the directional predictions for each O-series
outcome, along with the estimated effect size derived from preliminary analysis using
2020 cycle data. These estimates will be refined in O.1--O.4.

\begin{table}[ht]
\centering
\caption{Quantitative Predictions: \textsc{bisect} vs.\ Enacted Gerrymanders}
\label{tab:predictions}
\begin{tabular}{llll}
\toprule
Dimension & Direction & Point Estimate & Paper \\
\midrule
Competitiveness & Bisect $>$ enacted & $+4.2$ pp more competitive seats & O.1 \\
Turnout & Bisect $>$ enacted & $+1.8$ pp higher VAP turnout & O.2 \\
Polarization & Bisect $<$ enacted & $-0.08$ DW-NOMINATE gap reduction & O.3 \\
Constituent distance & Bisect $<$ enacted & $-14$ min mean driving time & O.4 \\
\bottomrule
\end{tabular}
\begin{tablenotes}
\small
\item Note: Estimates are preliminary, based on single-seed \textsc{bisect} comparison
  to enacted maps in 12 states with documented partisan gerrymandering. Confidence
  intervals will be reported in O.1--O.4.
\end{tablenotes}
\end{table}

\subsection{Scope Conditions}

The compactness-outcomes hypothesis applies most cleanly to states with documented
partisan gerrymandering, where the enacted map is systematically less compact than
the algorithmic alternative. In states without partisan gerrymandering (e.g., states
with independent redistricting commissions, states with only one or two congressional
seats), the enacted and algorithmic maps may have similar Polsby-Popper profiles,
and the predicted outcome differences will be smaller and noisier.

The O-series therefore stratifies its analysis by gerrymandering status: states
classified as having partisan gerrymanders (based on $|\text{EG}| > 0.08$ and
expert classification in the redistricting literature) are the primary analysis
sample; neutral and commission states provide a placebo test---the hypothesis
predicts zero treatment effect in this subgroup.
